\section{Theoretical Background}


\subsection{The 2D Ising Model}

The two-dimensional square-lattice Ising model is a paradigmatic system for studying phase transitions and critical phenomena. Spins $s_i = \pm 1$ occupy the sites of the $L \times L$ lattice with periodic boundary conditions (PBC) and interact with their nearest neighbors. The Hamiltonian for the system is given by:

\begin{equation}
H = -J \sum_{\langle i,j \rangle} s_i s_j - h \sum_i s_i,
\end{equation}

where $J$ is the interaction strength, $h$ is an external magnetic field, and the sum $\langle i,j \rangle$ runs over nearest neighbor pairs.

In zero external field ($h = 0$), the model exhibits a second-order phase transition at a critical temperature $T_c$, separating a low-temperature ferromagnetic phase (with non-zero spontaneous magnetization) from a high-temperature paramagnetic phase (with zero magnetization). For the 2D Ising model, the critical temperature is known exactly from Onsager's solution:
\begin{equation}
    \frac{k_B T_c}{J} = \frac{2}{\ln(1 + \sqrt{2})} \approx 2.2691853,
\end{equation}
or equivalently,
\begin{equation}
    J \beta_c = J \frac{1}{k_B T_c} \approx 0.4406868,
\end{equation}
where $k_B$ is the Boltzmann constant. \cite{2DIsingCritTemp}

The order parameter is the magnetization per spin, defined as:
\begin{equation}   
    m = \frac{1}{N} \sum_{i=1}^{N} s_i,
\end{equation}
where $N = L^2$ is the total number of spins in the lattice.

The zero-field susceptibility $\chi$ is defined thermodynamically as:
\begin{equation}
    \chi = \frac{\partial \langle m \rangle}{\partial h} \Big|_{h=0}.
\end{equation}

By the fluctuation-dissipation theorem (FDT), this equals to
\begin{equation}
    \chi = \beta N (\langle m^2 \rangle - \langle m \rangle^2).
\end{equation}

In finite simulations at $h = 0$, the  $\mathbb{Z}_2$ symmetry implies $\langle m \rangle = 0$ unless symmetry is explicitly broken. For thermodynamic estimates we use $m$ in the fluctuation formula for $\chi$, but for visualization we plot $\langle |m| \rangle$ in figures (in particular for the Wolff algorithm), since $|m|$ avoids sign cancellations and makes the ordered phase evident.

In practice, this derivative is often approximated by a finite-difference scheme using a small nonzero field $h_0$:
\begin{equation}
\chi \approx \frac{\langle M \rangle_{+h_0} - \langle M \rangle_{-h_0}}{2 h_0}.
\end{equation}


\subsection{Monte Carlo Dynamics}

\subsubsection{Markov Chains and Master Equations}

Monte Carlo updates define a Markov chain on the configuration space with a transition operator that has the Boltzmann distribution as its stationary measure.
For relaxational models, such as the stochastic Ising model, the time-dependent behavior is described by a master equation (\cite{LandauMonteCarlo})
\begin{equation} \label{equ:MasterEqu}
\frac{\partial P_n(t)}{\partial t}
= - \sum_{n \neq m} \left[ P_n(t) W_{n \to m} - P_m(t) W_{m \to n} \right],
\end{equation}
where $P_n(t)$ is the probability of the system being in state $n$ at time $t$, and $W_{n \to m}$ is the transition rate for $n \to m$.

The solution to the master equation is a sequence of states, but the time variable is a stochastic quantity which does not represent true time.
A relaxation function $\phi(t)$ can be defined which describes time correlations \emph{within equilibrium}
\begin{equation}
\phi_{MM}(t)
= \frac{\langle M(0)M(t)\rangle - \langle M \rangle^2}
{\langle M^2 \rangle - \langle M \rangle^2}.
\end{equation}

When normalized in this way, the relaxation function is $1$ at $t = 0$ and decays to zero as $t \to \infty$. It is important to remember that for a system in equilibrium any time in the sequence of states may be chosen as the `$t=0$' state. The asymptotic, long time behavior of the relaxation function is exponential, i.e.
\begin{equation}
\phi(t) \to e^{-t/\tau},
\end{equation}
where the correlation time $\tau$ diverges as $T_c$ is approached.


\subsubsection{Dynamic Critical Phenomena}
Following \cite{LandauMonteCarlo}, the dynamic (relaxational) critical behavior can also be expressed in terms of a power law
\begin{equation}
\tau \propto \xi^{z} \propto \varepsilon^{-\nu z},
\end{equation}
where $\xi$ is the (divergent) correlation length, $\varepsilon = |1 - T/T_c|$, and $z$ is the dynamic critical exponent.

In a finite system of linear size $L$, the divergence of the correlation
length at criticality is limited by the system size.
Exactly at $T_c$, the correlation length therefore saturates at $\xi \sim L$. Inserting this finite-size cutoff into the critical slowing-down relation $\tau \propto \xi^{z}$ gives
\begin{equation}
\tau_c \sim L^{z}.
\end{equation}
Thus, at the critical temperature, the relaxation time is no longer
controlled by the reduced temperature $\varepsilon$, but instead by the
system size itself.
As $L$ increases, correlations require more time to propagate across the entire system, leading to a relaxation time that scales as a power law in $L$ with exponent $z$.

A second relaxation time, the integrated relaxation time, is defined by the integral of the relaxation function
\begin{equation}
\tau_{\mathrm{int}} = \int_{0}^{\infty} \phi(t)\, dt.
\end{equation}
This quantity is of particular importance for the determination of errors and is expected to diverge with the same dynamic exponent as the `exponential' relaxation time, as both are governed by the same long-time behavior of $\phi(t)$.

